\section{Related Work}
\label{sec:related_work}

Our work lies at the intersection of parameter-efficient transfer learning, weight-space model merging, dynamic activation-space ensembling, and statistical learning theory. 

\subsection{Parameter-Efficient Fine-Tuning and Multi-Tenant Serving}
Parameter-Efficient Fine-Tuning (PEFT) techniques, such as bottleneck adapters \cite{houlsby2019parameter}, prefix tuning \cite{li2021prefix}, prompt tuning \cite{lester2021power}, and Low-Rank Adaptation (LoRA) \cite{hu2021lora}, enable efficient adaptation of large foundation models to diverse downstream tasks. Serving multiple concurrent adapters at scale has motivated systems-level frameworks like S-LoRA \cite{sheng2023slora} and Punica \cite{chen2023punica}, which optimize memory allocation and batching configurations for multi-tenant deployments. Additionally, LoraHub \cite{huang2024lorahub} proposes dynamic composition of multiple LoRA adapters via derivative-free optimization. However, these systems assume that task identities are known a priori or focus on optimizing homogeneous batching, rather than dynamically routing heterogeneous streams where task labels are unknown at test time.

\subsection{Weight-Space Model Merging}
Model merging aims to unify multiple specialized networks into a single cohesive model without incurring the cost of joint multi-task retraining. Linear weight averaging (Model Soups) \cite{wortsman2022model} works well for models fine-tuned from a shared initialization on similar tasks. For heterogeneous tasks, more advanced approaches have been developed. Task Arithmetic \cite{ilharco2022editing} computes task vectors by subtracting base weights from fine-tuned weights and adds them linearly. TIES-Merging \cite{yadav2023ties} resolves interference between task vectors by pruning small values, determining dominant signs, and scaling. DARE \cite{yu2023dare} randomly drops weights and rescales them to preserve the performance of the base model. To merge models with disparate architectures or permutations, Git Re-Basin \cite{ainsworth2022git} and ZipIt! \cite{stoica2023zipit} align hidden representation spaces before merging. Matena \& Raffel \cite{matena2021merging} propose merging models using Fisher Information. Although computationally efficient and parameter-free at serving time, weight-space merging suffers from \textit{heterogeneity collapse} when evaluated on mixed-task batches because a single static set of weights cannot adapt on-the-fly to different inputs.

\subsection{Dynamic Model Merging and Routing}
Dynamic model merging resolves heterogeneity collapse by routing inputs sample-wise to task-specific parameters at run-time. This is conceptually related to Mixture-of-Experts (MoE) models \cite{shazeer2017outrageously, fedus2022switch, riquelme2021scaling}, which utilize learned routing networks to dispatch tokens or inputs to specialized sub-networks. However, training MoEs from scratch is extremely expensive and susceptible to routing collapse. 
For pre-trained models on the edge, Parameter-Free Subspace Routing (PFSR) \cite{pfsr2025} uses classification heads to determine task membership, and Micro-Batch Homogenization (MBH) \cite{mbh2025} groups similar queries to optimize throughput. However, running sequential forward passes for different tasks scales as $O(K)$, creating prohibitive latency. 
Single-Pass Activation-Space Dynamic Blending (SPS) \cite{chatterjee2024robustness} solves this by performing dynamic activation-space ensembling inside a single forward pass, achieving $O(1)$ scaling. Zero-Shot Centroid Alignment (ZCA) uses cosine similarity against early-layer centroids to route without the "late-stage routing paradox" (running the backbone twice). Nonetheless, existing SPS and ZCA implementations (e.g., SABLE) rely on empirical heuristics, rendering them highly sensitive to heteroscedastic domain noise and representation anisotropy.

\subsection{PAC-Bayesian Learning Theory}
The PAC-Bayesian framework, originally introduced by McAllester \cite{mcallester1999pac, mcallester2003pac}, provides a principled learning-theoretic foundation for analyzing randomized classifiers (Gibbs policies) \cite{catoni2007pac}. Unlike classical PAC bounds, PAC-Bayesian bounds depend on a prior distribution over the hypothesis space and a posterior distribution learned from data, bounding the out-of-sample risk using the Kullback-Leibler (KL) divergence between them. PAC-Bayes has been successfully applied to analyze generalization in deep networks \cite{guedj2019primer, alquier2021user}. In this paper, we leverage the PAC-Bayesian framework not merely as an offline analytical tool, but as an \textit{active optimization objective} inside the model-merging router, providing provable generalization guarantees for dynamic on-device model serving.
